Instruction-conditioned deception in large language models is a practical safety problem, but many prior detectors rely on internal activations that are unavailable in black-box deployments. We ask a simpler question: can output text alone reveal when a model is pushed to lie? We run paired generation with \gptmodel on 120 \truthfulqa questions, producing truthful and lie-roleplay answers under matched decoding settings. We then test separability with both a text classifier and two-sample distribution tests.

Our output-only classifier (\tfidflogreg) reaches an AUROC of 0.883 on a held-out question split, with permutation-test $p=9.99\times10^{-4}$ and bootstrap 95\% CI $[0.798, 0.952]$. Independent distribution testing also rejects equal-distribution null hypotheses: \mmd is 0.02246 with permutation $p=0.001996$. The signal is robust to seed shift, where AUROC increases to 0.910 on unseen seed-7 generations. Feature-level analysis finds one strong univariate marker (punctuation ratio, FDR-corrected $p=1.97\times10^{-13}$, Cohen's $d=0.83$), while most predictive signal remains multivariate.

These results establish a reproducible black-box baseline for deception-risk monitoring from text alone. The findings suggest that lie-inducing instructions create stable output-distribution shifts that can be detected without privileged model access.
